\section{Limitaciones y trabajo pendiente}
\label{sec:limitaciones}

\subsection{Limitaciones de lo verificado}

\begin{itemize}
  \item \textbf{La placa no está fabricada.} Nada eléctrico ni de radiofrecuencia se ha medido: ni el consumo en sueño, ni el rizado del riel de 3.3~V al transmitir, ni la línea de 50~$\Omega$, ni el fix del GNSS con el parche, ni el alcance del LoRa, ni la enumeración por USB.
  \item \textbf{El DRC no es una revisión completa.} Comprueba las reglas cargadas. No sustituye a la revisión de integridad de señal ni a la de radiofrecuencia, y no se auditó que el layout cumpla los criterios de zonificación de la sección~4.
  \item \textbf{Los datos de la \pBoard{} son indicativos.} Otro procesador, otro receptor con menos constelaciones y otra antena. El resultado de la PoC en cielo parcial (\pCepLow{} a \pCepHigh~m de CEP50) no predice el de la PCB.
  \item \textbf{El presupuesto de consumo es calculado} con valores típicos de datasheet, con consumos todavía por sumar. El margen de \cMargin~\textmu A frente al límite de RNF-03 es estrecho.
\end{itemize}

\subsection{Hallazgos de las revisiones independientes}

Los dos hallazgos MAJOR de la revisión de RF tienen una decisión tomada el 2026-11-26; los de la revisión del BOM siguen sin decisión.
El texto completo está en \ruta{docs/design-reviews/}.
\begin{itemize}
  \item RF-01 (revisión de RF), \emph{resuelto en los documentos}: el perfil LoRa pasa a canales de 500~kHz dentro de 915.25 a 927.75~MHz. El firmware del LoRa aún no existe, y falta la confirmación de la ANE antes de cualquier prueba de campo.
  \item RF-02 (revisión de RF), \emph{decisión tomada, no aplicada}: la entrada del GNSS puede superar su máximo absoluto por los transmisores del propio collar. Esta PCB no lleva filtro y no se modifica; la siguiente revisión añade una huella de filtro SAW con puente de 0~$\Omega$. Mientras tanto, el firmware no transmite por LoRa ni por Wi-Fi con el GNSS activo, y el aislamiento se mide en la primera placa.
  \item BOM-02 y BOM-03 (revisión del BOM), \emph{abiertos}: el proyecto de Altium tiene parámetros de fabricante y proveedor equivocados en 25 piezas y no marca las piezas que no se montan (DNP), así que su lista de materiales no sirve para pedir piezas tal cual.
  \item BOM-04 (revisión del BOM), \emph{abierto}: el módulo ESP32-S3 admite 65~\textdegree C de ambiente, igual al máximo de diseño, y una carcasa sellada al sol lo supera.
\end{itemize}

\subsection{Documentos por corregir}

\begin{itemize}
  \item \ruta{docs/ARCHITECTURE.md}, secciones 2.2 y 2.3, aún da el divisor del MPPT como 402~k$\Omega$ / 100~k$\Omega$ y $R_{CS}$ como 0.12~$\Omega$. El esquemático tiene R8 + R9 = 400~k$\Omega$ con R10 = 100~k$\Omega$ (\cVmppt~V) y R12 = 0.47~$\Omega$ (\cIcs~mA). Esta nota sigue el esquemático.
  \item \ruta{docs/REQUIREMENTS.md}, RF-03, nombra el zumbador PKMCS0909E4000-R1, y el esquemático lleva el PS1240P02BT (BZ1).
\end{itemize}

\subsection{Siguientes pasos}

\begin{enumerate}
  \item Exportar el reporte ERC de Altium (no se exportó) y añadirlo a la tabla de verificación.
  \item Decidir los hallazgos BOM-02 a BOM-04, aplicar la huella del filtro GNSS en la siguiente revisión de la PCB y confirmar el perfil LoRa con la ANE.
  \item Fabricar la placa y medir: enumeración y programación por USB, fix del GNSS con el parche, rizado del riel de 3.3~V en transmisión, alcance del LoRa y consumo en sueño. Con ello, la tabla de cumplimiento pasa de ``por diseño'' a ``medido en la PCB''.
  \item Portar el firmware al ESP32-S3: backend UART del MAX-M10S, LoRa, IMU y sueño profundo.
\end{enumerate}
